\section{Introduction}

Modern cloud environments provide hundreds of distributed systems and
services that are readily available for developers to build a
large-scale application. System architects and developers following the
microservice or service-oriented designs can freely
choose and compose the building blocks in a modular fashion and plug them
into the application to use.

However, different distributed systems provide different consistency semantics
and the mismatch of the chosen system's semantics and the
application's expected semantics could introduce subtle silent
semantic bugs~\cite{silentbugs}. For example, a user may connect a key-value store
observing eventual consistency to an application that expects a
linearizable storage service. In this case, the application could
read a non-up-to-date copy of the data and application could produce
an erroneous outcome.

The user could carefully review the application needs and choose a
matching system, but when the application depends on multiple
distributed services, checking the overall compatibility between the
application logic and multiple systems easily gets out of hand. Even
if one successfully finds the right combination, for example, through
extensive testing, the user could change the application code or
plug in different systems due to management, cost, and performance
reasons. Then, the sophisticated compatibility check needs to be done again.

% When you have fixed application logic, you may choose two
% compatible systems. Two systems may collectively satisfy the
% application requirement, and changing one may invalidate both.
% NOTE: addressed below

Formal verification can be another way to address such challenges,
where the correctness of the code can mathematically guarantee the
compatibility between the code and the system~\cite{lampson2021hints}. 
However, depending on
the complexity of the code, which could include non-essential parts
for reasoning about compatibility, the formal verification process
can take extensive time and effort~\cite{certikos-osdi16, klein2009sel4, wormspace}. 
Furthermore, depending on the
formal verification techniques, the amount of proofs that needs to be
written could be several times compare to the lines of code and small
change could easily break the entire proof~\cite{klein2009sel4}. 
Thus, formal verification
leads to a similar problem to testing where the compatibility check
cannot agilely keep up with the changes in underlying systems and 
applications. \todo{Lindsey's work on trace based analysis??}

In this paper, we propose \sysname{}, a light-weight \sysnamelong{}. \sysname{}
helps developers to quickly encode their application logic and
requirements and check the compatibility with underlying distributed
systems using an SMT-based approach. 
%to encode and analyze 
%the safety properties of different consistency semantics and
%their compositions.

\sysname{} includes axiomatic models of distributed consistency
semantics, including eventual consistency,
linearizability~\cite{viotti2016consistency}, four session
guarantees~\cite{terry1994session}, pipeline RAM (PRAM) consistency~\cite{pram},
and causal consistency~\cite{causaldef}, which are commonly used in both
industry and academia~\needcite{}. Users can use the models as
building blocks to test the compatibility of these models with their
applications. Users can add additional models and test the
correctness and compatibility with existing models as well.

Although axiomatic definitions of distributed consistency semantics
exist in literatures~\needcite{}, encoding the details into the SMT 
clauses for compatibility checks are challenging as the definitions
glosses over a few key details. For example, the commonly used concept
of \textit{visibility} for defining weak distributed consistency 
semantics is subtle, but to our knowledge, no literature provides
a formal definition. \sysname{} fills in these missing gaps to realize
full-fledged formal models for the compatibility check. 

Each functionality in the application logic in \sysname{} is encoded
as a directed acyclic graph (DAG) in a systematic way where the entire 
application can be encoded into a directed multi graph format which we call the \sysgraph{}. 
Each node in the \sysgraph{} represents key functionality in the 
application which requires certain properties from lower-level nodes 
it depends on and provides certain properties to the upper-level nodes.
The edges are used to specify more specific requirements and condition that each
functionality needs to enforce in accessing the lower-level nodes.
The leaf nodes are typically the external distributed service that
follows one of the pre-built distributed consistency semantic models
in \sysname{}.

The \sysgraph{} model is expressive enough to model an application logic 
from a high-level and trace complex dependencies that cross multiple
application logics and distributed systems that the application
replies on. \sysname{} can trivially check simple application that
relies on one distributed service and can span to a complex
micro-service application that leverages multiple modules and
distributed systems.

Once the user specifies the application logic, \sysname{} can quickly
check the compatibility with the plugged in distributed system. It
could be testing what the user chose, or alternatively, \sysname{}
can explore different combinations of existing models to
find distributed systems with compatible consistency semantics. These
approaches enable fast exploration and iteration of alternative
system designs, guiding the development to adhere to intended
consistency guarantees from the very beginning.

%%%%%% Examples %%%%%%    % NOTE: addressed here
%%%%%% Examples %%%%%%    {
%%%%%% Examples %%%%%%      \color{red}
%%%%%% Examples %%%%%%      Consider a Not So Simple Storage Service (NS4) that advertises
%%%%%% Examples %%%%%%      itself to provide linearizability guarantee to users' read/write to objects.
%%%%%% Examples %%%%%%      NS4's architecture is based on an API layer that provides
%%%%%% Examples %%%%%%      arbitration between users and the
%%%%%% Examples %%%%%%      underlying monolithic storage layer
%%%%%% Examples %%%%%%      with read-your-writes guarantee. In this case, \sysname{} can be used to
%%%%%% Examples %%%%%%      model the abstract semantical properties of NS4, and check whether the
%%%%%% Examples %%%%%%      modeled NS4 have advertised user-facing linearizable guarantee.
%%%%%% Examples %%%%%%      Three nodes (client, API, and storage) and two edges (client to
%%%%%% Examples %%%%%%      API, API to storage) are
%%%%%% Examples %%%%%%      added to the DAG, where the client node requires linearizability
%%%%%% Examples %%%%%%      from the API node,
%%%%%% Examples %%%%%%      API node exports linearizability to the client, and API node
%%%%%% Examples %%%%%%      requires read-your-writes from the storage node.
%%%%%% Examples %%%%%%      The storage node exports read-your-writes to the API node. With
%%%%%% Examples %%%%%%      this modeling, \sysname{}
%%%%%% Examples %%%%%%      will be able to find a satisfying assignment for the DAG where
%%%%%% Examples %%%%%%      user-facing guarantee is linearizability.
%%%%%% Examples %%%%%%      However, if we swap the API layer requirement from the storage
%%%%%% Examples %%%%%%      layer from read-your-writes to
%%%%%% Examples %%%%%%      to PRAM consistency, \sysname{} will be able to find that the leaf
%%%%%% Examples %%%%%%      level semantical requirement can no longer be satisfied.
%%%%%% Examples %%%%%%    }

%%%%%% %%%%%%    \jiyong{Some challenges in doing this?}
%%%%%% %%%%%%    \yifei{What aspect should we talk about? Preliminary encoding on
%%%%%% %%%%%%      semantics? Debugging? Unresolved encoding on models with "eventual"
%%%%%% %%%%%%    visibility (see eventual source code)?}



As case studies we use \sysname{} to model two microservice
applications~\cite{deathstar}, e-commerce service and movie streaming
service, and a causally consistent storage service built on top of
multiple weakly consistent distributed storage systems~\cite{boltoncc} to
check the compatibility with underlying distributed systems and
explore other compatible semantics. Our framework is capable of
checking the compatibility under XXX seconds and explore other
compatible systems in XXX seconds.

%%%%%% %%%%%%    \jiyong{Some insights we gained?}
%%%%%% %%%%%%    \yifei{
%%%%%% %%%%%%      This is a project with a lot of moving parts, modeling a system,
%%%%%% %%%%%%      like in real life, have multiple ways to achieve the same goal.
%%%%%% %%%%%%      The level of abstraction we've chosen made it not too hard to check
%%%%%% %%%%%%      real system's behavior but that part is not implemented.
%%%%%% %%%%%%      Personally I feel like some part of the project would be better
%%%%%% %%%%%%      written in a proper theorem prover (representation of ops, binary
%%%%%% %%%%%%      relations), and discharge part of semantic constraints to SMT solvers.
%%%%%% %%%%%%    }

This paper makes the following contributions:
\itemparen{
\item We present \sysname{}, the first compatibility checking tool for
  application that is composed of multiple distributed services.
\item We provide the missing formal definitions to fully express
  and check distributed weak consistency semantics. 
\item We present a systematic way of encoding complex applications
  using the \sysgraph{} structure and checking their compatibility
  collectively using the SMT-based approach.
\item We carry out a case study on three representative distributed
  applications to demonstrate the workings of \sysname{}.
}

%%%%    % FIXME: full rewrite
%%%%    {\color{red}
%%%%
%%%%      Modern distributed systems often relax strict consistency
%%%%      guarantees to achieve
%%%%      scalability, fault tolerance, and low latency operations.
% Understanding and
%%%%      comparing different forms of relaxed consistency semantics
% is crucial for both
%%%%      system designers and researchers. Yet, systematically
% modeling and reasoning
%%%%      about these semantics, as well as their compositional
% properties, remained a
%%%%      understudied area.
%%%%
%%%%      Our framework aims to bridge this gap by providing a
% SMT-based approach to
%%%%      encode and analyze the safety properties of weaker consistency
%%%%      models (comparing
%%%%      with linearizability \cite{herlihy1990linearizability}), and their
%%%%      compositions under
%%%%      the assumption of livenss properties are always satisfied,
% i.e. eventual
%%%%      consistency \cite{petersen1997flexible} is guaranteed (replicas
%%%%      eventually converge),
%%%%      and we are layering everything on top of it. By offering a per-system
%%%%      (differentiated by symbols), logical representation of
% operations, our system
%%%%      can differentiates operations belonging to distinct systems
% and/or components
%%%%      through the use of unique symbols (see
% tests/test\_cross\_causal.py), enabling
%%%%      analysis of cross-system interactions. Additionally, solver
% backends can
%%%%      discharge proof witnesses on the satisfiability of specific
% consistency
%%%%      properties, or the compatibility and composability between
% multiple models.
%%%%      These approaches enable fast exploration and iteration of
% alternative system
%%%%      designs, guiding the development adhere to intended consistency
%%%%      guarantees from
%%%%      the very beginning.
%%%%
%%%%    }
